\documentclass[11pt, reqno]{article}

\usepackage{amsmath, amsthm, amssymb}
\usepackage{enumitem}
\usepackage{tcolorbox}
\usepackage{hyperref}
\usepackage{tikz}
\usepackage{tikz-cd}
\usepackage{pgfplots}
\pgfplotsset{compat=1.18}
\usetikzlibrary{arrows.meta}
\usepackage{mathrsfs}
\usepackage{fancyhdr}
\usepackage[bottom=0.75in, top=1in, left=0.5in, right=0.5in]{geometry}
\usepackage{array}   % for \newcolumntype macro
\newcolumntype{L}{>{$}l<{$}}

\theoremstyle{plain}
\newtheorem*{theorem}{Theorem}
\newtheorem*{proposition}{Proposition}
\newtheorem{exercise}{Exercise}
\newtheorem*{lemma}{Lemma}
\newtheorem*{corollary}{Corollary}

\theoremstyle{definition}
\newtheorem*{definition}{Definition}
\newtheorem*{example}{Example}

\theoremstyle{remark}
\newtheorem*{remark}{Remark}

\renewcommand{\phi}{\varphi}
\renewcommand{\epsilon}{\varepsilon}
\renewcommand{\emptyset}{\varnothing}

\newcommand{\RR}{\mathbb{R}}
\newcommand{\ZZ}{\mathbb{Z}}
\newcommand{\NN}{\mathbb{N}}
\newcommand{\CC}{\mathbb{C}}
\newcommand{\QQ}{\mathbb{Q}}

\DeclareMathOperator{\ima}{\text{im}}

\begin{document}

\topmargin=-40pt
\rhead{Henry Woodburn}
\lhead{Summer 2026}
\renewcommand{\headrulewidth}{1pt}
\renewcommand{\headsep}{20pt}
\thispagestyle{fancy}

{\Huge \bfseries \noindent K-Homology Notes}

\section{Chapter 1}

Some Definitions: We say that a representation of a $C^*$ algebra
\[
    \pi: A \to B(H)
\]
is:
\begin{enumerate}
    \item Faithful: $\pi$ is injective
    \item Nondegenerate: $\pi(1) = 1$
    \item Ample: No operator of $A$ acts as a compact operator.
\end{enumerate}

\textbf{Setting}: Let $H$ be a complex Hilbert space equipped with a hermitian inner 
product. We denote by $B(H)$ the subalgebra of bounded operators. The \textit{adjoint}
$T^*$ of an operator $T$ is defined by $\langle Tv, w\rangle = \langle v, T^* w\rangle$
for all $v, w$. The adjoint is an involution. Here we also have the $C^*$ identity
\[
    \|T^*T\| = \|T\|^2.
\]
\begin{example}
    Any $*$ subalgebra of $B(H)$ is a $C^*$ algebra. Conversely,
    a Banach-$*$ algebra with the $C^*$ identity is called a $C^*$ algebra. 
\end{example}

\begin{theorem}[Gelfand]
    If $T$ is normal, then 
    \[
        \phi: C^*(T) \to C(\text{Spec}(T))
    \]
    is an isomorphism.
\end{theorem}

\begin{definition}
    The \textit{dual} of a Banach algebra $A$ is defined to be 
    \[
        \hat{A} = \{f: A \to \CC: f\ \text{a nonzero continuous homomorphism}\}.
    \]
\end{definition}

\section{Meeting 2}

\textbf{Some Notes by Chenxi:}

\begin{enumerate}
    \item[1.] Functional calculus can be viewed as $C(X)$ acting on $L^2(X, \mu)$.
    
    \textbf{Fact:} We can write 
    \[
        H = \sum_{v} \text{span}(v, Tv, \dots)
    \]
    where the $v$ are obtained using Zorn's lemma. 

    \textbf{Fact:} We have $\langle P(T) v, v\rangle$ for any polynomial $P$, 
    and we use Riesz to extend to arbitrary continuous $f$.

    Then we have a functional on $H$, so there is a Borel measure $\mu$ such that 
    \[
        \langle f(T)v, v\rangle = \int_X f(x) d\mu
    \]

    \textbf{Unbounded things:} Extend using mobius transform. 
\end{enumerate}

\textbf{Notes by Lily:}

\begin{definition}
    We define the \textbf{Caulkin algebra} to be the quotient $B(H)/K(H)$, where 
    $B$ and $K$ denote the spaces of bounded or compact operators.

    The \textbf{Fredholm Index}, $\text{Ind}(T) = \dim \ker(T) - \dim \text{coker}(T)$.
\end{definition}

\begin{theorem}[Atkinson]
    $T \in B(H)$ is fredholm iff its image in the Calkin algebra is invertible. 
\end{theorem}

We also have $\text{Ind}(ST) = \text{Ind}(S) + \text{Ind}(T)$.

\begin{definition}
    Let $T \in B(H)$. The \textbf{Essential Spectrum} of $T$ is defined as the ordinary
    spectrum of the image of $T$ in the Calkin algebra. 

    $T_1$ and $T_2$ on $H_1$ and $H_2$ are essentially unitarily equivalent if 
    there is a unitary map $\mathcal{U}: H_1 \to H_2$ st $T_1 - \mathcal{U}^* T \mathcal{U}$ 
    is compact. 
\end{definition}

\begin{theorem}[Weyl-von Neumann]
    Every bounded self adjoint operator on a separable Hilbert space is an 
    arbirarily small compact perturbation of a diagonal element.
\end{theorem}

\end{document}